% !TeX root = part01.tex

\documentclass[../final.tex]{subfiles}

\begin{document}

% ============================================================
% Семинар 1
% ============================================================

\practice{1}{05.09.2026}{Кристаллография}


% ============================================================
% Задача 1
% ============================================================

\section{Задача 1. ОЦК как решётка Браве}

\textbf{Условие.}
Показать, что объёмно-центрированная кубическая решётка
является решёткой Браве.


\begin{rbox}[left=18pt,right=18pt,top=14pt,bottom=10pt]{[]}
    \centering

    \includegraphics[
        width=0.88\linewidth
    ]{image/practice_01/bcc_lattice.png}

    \captionsetup{
        font=footnotesize,
        labelfont=bf,
        skip=5pt
    }

    \captionof{figure}{
        Объёмно-центрированная кубическая решётка.
        Слева показана обычная кубическая ячейка с узлами
        в вершинах и одним узлом в центре.
        Справа изображены три соседние ячейки,
        имеющие общий узел; векторы
        $\mathbf a_1$, $\mathbf a_2$, $\mathbf a_3$
        направлены из общего узла к объёмно-центрированным узлам
        соседних ячеек.
        Также показаны координатные оси $x$, $y$, $z$.
    }

    \label{fig:bcc_lattice}
\end{rbox}


Рассмотрим объёмно-центрированную кубическую решётку
с параметром обычной кубической ячейки $a$.

Для того чтобы показать, что это решётка Браве, достаточно
подобрать три примитивных вектора трансляции так, чтобы любой
узел решётки можно было получить их целочисленной линейной
комбинацией.

Выберем
%
\begin{equation*}
    \mathbf a_1
    =
    \frac{a}{2}
    \left(
        \mathbf e_x
        +
        \mathbf e_y
        -
        \mathbf e_z
    \right),
\end{equation*}
%
\begin{equation*}
    \mathbf a_2
    =
    \frac{a}{2}
    \left(
        -\mathbf e_x
        +
        \mathbf e_y
        +
        \mathbf e_z
    \right),
\end{equation*}
%
\begin{equation*}
    \mathbf a_3
    =
    \frac{a}{2}
    \left(
        \mathbf e_x
        -
        \mathbf e_y
        +
        \mathbf e_z
    \right).
\end{equation*}

Эти векторы направлены из одного узла ОЦК
в три ближайших соседних узла.

Тогда положение любого узла решётки можно записать в виде
%
\begin{equation*}
    \boxed{
        \mathbf R
        =
        n_1\mathbf a_1
        +
        n_2\mathbf a_2
        +
        n_3\mathbf a_3
    },
\end{equation*}
%
где
%
\begin{equation*}
    n_1,n_2,n_3\in\mathbb Z.
\end{equation*}

Таким образом, вся решётка получается трансляциями
одного и того же узла.

Следовательно,
%
\begin{equation*}
    \boxed{
        \text{ОЦК является решёткой Браве.}
    }
\end{equation*}

\textbf{Комментарий.}
В обычной кубической ячейке ОЦК кажется, что имеются узлы
двух типов: вершины куба и центр куба. Однако это связано
только с выбором непримитивной кубической ячейки.

Если перейти к примитивной ячейке с векторами
$\mathbf a_1$, $\mathbf a_2$, $\mathbf a_3$, все узлы
оказываются трансляционно эквивалентными.


% ============================================================
% Задача 2
% ============================================================

\section{Задача 2. ГЦК как решётка Браве}

\textbf{Условие.}
Показать, что гранецентрированная кубическая решётка
является решёткой Браве.


% ------------------------------------------------------------
% РИСУНОК 2
%
% Вставить здесь рисунок ГЦК из рукописного конспекта:
% куб с узлами в вершинах и в центрах граней.
%
% На рисунке желательно показать:
%   1) оси x, y, z;
%   2) три примитивных вектора
%      \mathbf a_1, \mathbf a_2, \mathbf a_3;
%   3) центры граней, в которые приходят эти векторы.
% ------------------------------------------------------------


Рассмотрим гранецентрированную кубическую решётку.

В качестве примитивных векторов выберем
%
\begin{equation*}
    \mathbf a_1
    =
    \frac{a}{2}
    \left(
        \mathbf e_x+\mathbf e_y
    \right),
\end{equation*}
%
\begin{equation*}
    \mathbf a_2
    =
    \frac{a}{2}
    \left(
        \mathbf e_y+\mathbf e_z
    \right),
\end{equation*}
%
\begin{equation*}
    \mathbf a_3
    =
    \frac{a}{2}
    \left(
        \mathbf e_z+\mathbf e_x
    \right).
\end{equation*}

Геометрически эти векторы направлены из одного узла
к центрам трёх соседних граней куба.

Любой узел ГЦК можно представить как
%
\begin{equation*}
    \boxed{
        \mathbf R
        =
        n_1\mathbf a_1
        +
        n_2\mathbf a_2
        +
        n_3\mathbf a_3
    },
    \qquad
    n_1,n_2,n_3\in\mathbb Z.
\end{equation*}

Следовательно, все узлы решётки получаются из одного узла
с помощью трансляций.

Поэтому
%
\begin{equation*}
    \boxed{
        \text{ГЦК является решёткой Браве.}
    }
\end{equation*}

\textbf{Комментарий.}
Как и в случае ОЦК, обычная кубическая ячейка ГЦК
не является примитивной.

В ней содержится несколько узлов решётки, но это не означает
наличия нескольких атомов в базисе. После выбора примитивных
векторов все узлы оказываются эквивалентными.


% ============================================================
% Обратная решётка
% ============================================================

\section*{Обратная решётка}


% ============================================================
% Задача 3
% ============================================================

\section{Задача 3. Обратная решётка к ОЦК}

\textbf{Условие.}
Показать, что обратная решётка к ОЦК является ГЦК.

Для прямой решётки с примитивными векторами
%
\begin{equation*}
    \mathbf a_1,
    \qquad
    \mathbf a_2,
    \qquad
    \mathbf a_3
\end{equation*}
%
введём примитивные векторы обратной решётки:
%
\begin{equation*}
    \boxed{
        \mathbf b_1
        =
        2\pi
        \frac{
            \mathbf a_2\times\mathbf a_3
        }{
            \mathbf a_1\cdot
            (\mathbf a_2\times\mathbf a_3)
        }
    },
\end{equation*}
%
\begin{equation*}
    \boxed{
        \mathbf b_2
        =
        2\pi
        \frac{
            \mathbf a_3\times\mathbf a_1
        }{
            \mathbf a_1\cdot
            (\mathbf a_2\times\mathbf a_3)
        }
    },
\end{equation*}
%
\begin{equation*}
    \boxed{
        \mathbf b_3
        =
        2\pi
        \frac{
            \mathbf a_1\times\mathbf a_2
        }{
            \mathbf a_1\cdot
            (\mathbf a_2\times\mathbf a_3)
        }
    }.
\end{equation*}

Для этих векторов выполняется условие
%
\begin{equation*}
    \boxed{
        \mathbf a_i\cdot\mathbf b_j
        =
        2\pi\delta_{ij}
    }.
\end{equation*}

То есть, например,
%
\begin{equation*}
    \mathbf a_1\cdot\mathbf b_1=2\pi,
    \qquad
    \mathbf a_1\cdot\mathbf b_2=0,
    \qquad
    \mathbf a_1\cdot\mathbf b_3=0,
\end{equation*}
%
и аналогично для остальных пар.

Для ОЦК имеем
%
\begin{equation*}
    \mathbf a_1
    =
    \frac{a}{2}
    \left(
        \mathbf e_x+\mathbf e_y-\mathbf e_z
    \right),
\end{equation*}
%
\begin{equation*}
    \mathbf a_2
    =
    \frac{a}{2}
    \left(
        -\mathbf e_x+\mathbf e_y+\mathbf e_z
    \right),
\end{equation*}
%
\begin{equation*}
    \mathbf a_3
    =
    \frac{a}{2}
    \left(
        \mathbf e_x-\mathbf e_y+\mathbf e_z
    \right).
\end{equation*}


% ------------------------------------------------------------
% Вычисление b_1
% ------------------------------------------------------------

\subsection*{Вычисление первого вектора обратной решётки}

Найдём сначала числитель для $\mathbf b_1$:
%
\begin{equation*}
    \mathbf a_2\times\mathbf a_3
    =
    \frac{a^2}{4}
    \left[
        (-\mathbf e_x+\mathbf e_y+\mathbf e_z)
        \times
        (\mathbf e_x-\mathbf e_y+\mathbf e_z)
    \right].
\end{equation*}

Раскроем векторное произведение по членам:
%
\begin{align*}
    (-\mathbf e_x)\times\mathbf e_x
    &=
    0,
    \\
    (-\mathbf e_x)\times(-\mathbf e_y)
    &=
    \mathbf e_x\times\mathbf e_y
    =
    \mathbf e_z,
    \\
    (-\mathbf e_x)\times\mathbf e_z
    &=
    +\mathbf e_y,
    \\
    \mathbf e_y\times\mathbf e_x
    &=
    -\mathbf e_z,
    \\
    \mathbf e_y\times(-\mathbf e_y)
    &=
    0,
    \\
    \mathbf e_y\times\mathbf e_z
    &=
    \mathbf e_x,
    \\
    \mathbf e_z\times\mathbf e_x
    &=
    \mathbf e_y,
    \\
    \mathbf e_z\times(-\mathbf e_y)
    &=
    +\mathbf e_x,
    \\
    \mathbf e_z\times\mathbf e_z
    &=
    0.
\end{align*}

Поэтому
%
\begin{align*}
    \mathbf a_2\times\mathbf a_3
    &=
    \frac{a^2}{4}
    \left(
        \mathbf e_z
        +
        \mathbf e_y
        -
        \mathbf e_z
        +
        \mathbf e_x
        +
        \mathbf e_y
        +
        \mathbf e_x
    \right)
    \\
    &=
    \frac{a^2}{4}
    \left(
        2\mathbf e_x+2\mathbf e_y
    \right).
\end{align*}

Следовательно,
%
\begin{equation*}
    \boxed{
        \mathbf a_2\times\mathbf a_3
        =
        \frac{a^2}{2}
        \left(
            \mathbf e_x+\mathbf e_y
        \right)
    }.
\end{equation*}

Теперь найдём объём примитивной ячейки:
%
\begin{align*}
    \mathbf a_1\cdot
    (\mathbf a_2\times\mathbf a_3)
    &=
    \frac{a}{2}
    \left(
        \mathbf e_x+\mathbf e_y-\mathbf e_z
    \right)
    \cdot
    \frac{a^2}{2}
    \left(
        \mathbf e_x+\mathbf e_y
    \right)
    \\
    &=
    \frac{a^3}{4}
    (1+1).
\end{align*}

Здесь использована ортонормированность декартова базиса:
%
\begin{equation*}
    \mathbf e_x\cdot\mathbf e_x
    =
    \mathbf e_y\cdot\mathbf e_y
    =
    \mathbf e_z\cdot\mathbf e_z
    =
    1,
\end{equation*}
%
а смешанные произведения равны нулю.

Поэтому
%
\begin{equation*}
    \boxed{
        \mathbf a_1\cdot
        (\mathbf a_2\times\mathbf a_3)
        =
        \frac{a^3}{2}
    }.
\end{equation*}

Подставляем полученные выражения в определение $\mathbf b_1$:
%
\begin{align*}
    \mathbf b_1
    &=
    2\pi
    \frac{
        \dfrac{a^2}{2}
        (\mathbf e_x+\mathbf e_y)
    }{
        \dfrac{a^3}{2}
    }
    \\
    &=
    \frac{2\pi}{a}
    (\mathbf e_x+\mathbf e_y).
\end{align*}

Таким образом,
%
\begin{equation*}
    \boxed{
        \mathbf b_1
        =
        \frac{2\pi}{a}
        \left(
            \mathbf e_x+\mathbf e_y
        \right)
    }.
\end{equation*}

Аналогично получаем
%
\begin{equation*}
    \boxed{
        \mathbf b_2
        =
        \frac{2\pi}{a}
        \left(
            \mathbf e_y+\mathbf e_z
        \right)
    },
\end{equation*}
%
\begin{equation*}
    \boxed{
        \mathbf b_3
        =
        \frac{2\pi}{a}
        \left(
            \mathbf e_x+\mathbf e_z
        \right)
    }.
\end{equation*}


% ------------------------------------------------------------
% Сравнение с ГЦК
% ------------------------------------------------------------

\subsection*{Сравнение с примитивными векторами ГЦК}

Для ГЦК с параметром обычной кубической ячейки $a'$
примитивные векторы имеют вид
%
\begin{equation*}
    \mathbf a'_1
    =
    \frac{a'}{2}
    (\mathbf e_x+\mathbf e_y),
\end{equation*}
%
\begin{equation*}
    \mathbf a'_2
    =
    \frac{a'}{2}
    (\mathbf e_y+\mathbf e_z),
\end{equation*}
%
\begin{equation*}
    \mathbf a'_3
    =
    \frac{a'}{2}
    (\mathbf e_x+\mathbf e_z).
\end{equation*}

Полученные векторы обратной решётки имеют точно такую же
структуру.

Сравнивая коэффициенты,
%
\begin{equation*}
    \frac{a'}{2}
    =
    \frac{2\pi}{a},
\end{equation*}
%
получаем
%
\begin{equation*}
    \boxed{
        a'
        =
        \frac{4\pi}{a}
    }.
\end{equation*}

Таким образом,
%
\begin{equation*}
    \boxed{
        \text{обратная решётка к ОЦК является ГЦК.}
    }
\end{equation*}

\textbf{Комментарий.}
Аналогично можно показать обратное утверждение:
обратная решётка к ГЦК является ОЦК.

Таким образом, ОЦК и ГЦК являются взаимно обратными
типами кубических решёток.


% ============================================================
% Задача 4
% ============================================================

\section{Задача 4. Вектор обратной решётки и плоскости прямой решётки}

\textbf{Условие.}
Показать, что каждый вектор обратной решётки перпендикулярен
некоторому семейству плоскостей прямой решётки.


% ------------------------------------------------------------
% РИСУНОК 3
%
% Вставить здесь рисунок из рукописного конспекта.
%
% На рисунке:
%   1) из начала координат проведены векторы
%      n_1\mathbf a_1,
%      n_2\mathbf a_2,
%      n_3\mathbf a_3;
%   2) через их концы проведена плоскость;
%   3) в этой плоскости показаны два вектора
%      \mathbf c_2 и \mathbf c_3.
% ------------------------------------------------------------


Пусть плоскость пересекает направления примитивных векторов
в точках
%
\begin{equation*}
    n_1\mathbf a_1,
    \qquad
    n_2\mathbf a_2,
    \qquad
    n_3\mathbf a_3.
\end{equation*}

Построим два независимых вектора, лежащих в этой плоскости:
%
\begin{equation*}
    \boxed{
        \mathbf c_2
        =
        n_2\mathbf a_2
        -
        n_1\mathbf a_1
    },
\end{equation*}
%
\begin{equation*}
    \boxed{
        \mathbf c_3
        =
        n_3\mathbf a_3
        -
        n_1\mathbf a_1
    }.
\end{equation*}

Вектор
%
\begin{equation*}
    \mathbf c_2\times\mathbf c_3
\end{equation*}
%
перпендикулярен обоим векторам $\mathbf c_2$ и $\mathbf c_3$,
а значит, является нормалью к рассматриваемой плоскости.

Раскроем векторное произведение:
%
\begin{align*}
    \mathbf c_2\times\mathbf c_3
    &=
    \left(
        n_2\mathbf a_2-n_1\mathbf a_1
    \right)
    \times
    \left(
        n_3\mathbf a_3-n_1\mathbf a_1
    \right)
    \\
    &=
    n_2n_3
    [\mathbf a_2\times\mathbf a_3]
    -
    n_1n_2
    [\mathbf a_2\times\mathbf a_1]
    -
    n_1n_3
    [\mathbf a_1\times\mathbf a_3]
    +
    n_1^2
    [\mathbf a_1\times\mathbf a_1].
\end{align*}

Последний член равен нулю:
%
\begin{equation*}
    [\mathbf a_1\times\mathbf a_1]=0.
\end{equation*}

Поэтому
%
\begin{equation*}
    \mathbf c_2\times\mathbf c_3
    =
    n_2n_3
    [\mathbf a_2\times\mathbf a_3]
    -
    n_1n_2
    [\mathbf a_2\times\mathbf a_1]
    -
    n_1n_3
    [\mathbf a_1\times\mathbf a_3].
\end{equation*}

Используем антисимметрию векторного произведения:
%
\begin{equation*}
    [\mathbf a_2\times\mathbf a_1]
    =
    -[\mathbf a_1\times\mathbf a_2],
\end{equation*}
%
\begin{equation*}
    [\mathbf a_1\times\mathbf a_3]
    =
    -[\mathbf a_3\times\mathbf a_1].
\end{equation*}

Тогда
%
\begin{equation*}
    \boxed{
        \mathbf c_2\times\mathbf c_3
        =
        n_2n_3
        [\mathbf a_2\times\mathbf a_3]
        +
        n_1n_2
        [\mathbf a_1\times\mathbf a_2]
        +
        n_1n_3
        [\mathbf a_3\times\mathbf a_1]
    }.
\end{equation*}

Вынесем общий множитель $n_1n_2n_3$:
%
\begin{align*}
    \mathbf c_2\times\mathbf c_3
    &=
    n_1n_2n_3
    \left[
        \frac{1}{n_1}
        [\mathbf a_2\times\mathbf a_3]
        +
        \frac{1}{n_2}
        [\mathbf a_3\times\mathbf a_1]
        +
        \frac{1}{n_3}
        [\mathbf a_1\times\mathbf a_2]
    \right].
\end{align*}

Введём объём примитивной ячейки
%
\begin{equation*}
    \Omega
    =
    \mathbf a_1\cdot
    [\mathbf a_2\times\mathbf a_3].
\end{equation*}

Из определения обратных векторов имеем
%
\begin{equation*}
    \mathbf b_1
    =
    \frac{2\pi}{\Omega}
    [\mathbf a_2\times\mathbf a_3],
\end{equation*}
%
\begin{equation*}
    \mathbf b_2
    =
    \frac{2\pi}{\Omega}
    [\mathbf a_3\times\mathbf a_1],
\end{equation*}
%
\begin{equation*}
    \mathbf b_3
    =
    \frac{2\pi}{\Omega}
    [\mathbf a_1\times\mathbf a_2].
\end{equation*}

Следовательно,
%
\begin{equation*}
    [\mathbf a_2\times\mathbf a_3]
    =
    \frac{\Omega}{2\pi}\mathbf b_1,
\end{equation*}
%
\begin{equation*}
    [\mathbf a_3\times\mathbf a_1]
    =
    \frac{\Omega}{2\pi}\mathbf b_2,
\end{equation*}
%
\begin{equation*}
    [\mathbf a_1\times\mathbf a_2]
    =
    \frac{\Omega}{2\pi}\mathbf b_3.
\end{equation*}

Подставляем:
%
\begin{equation*}
    \mathbf c_2\times\mathbf c_3
    =
    \frac{\Omega}{2\pi}
    n_1n_2n_3
    \left(
        \frac{\mathbf b_1}{n_1}
        +
        \frac{\mathbf b_2}{n_2}
        +
        \frac{\mathbf b_3}{n_3}
    \right).
\end{equation*}


% ------------------------------------------------------------
% Индексы Миллера
% ------------------------------------------------------------

\subsection*{Связь с индексами Миллера}

Теперь воспользуемся определением индексов Миллера.

Если плоскость пересекает направления примитивных векторов
в точках
%
\begin{equation*}
    n_1\mathbf a_1,
    \qquad
    n_2\mathbf a_2,
    \qquad
    n_3\mathbf a_3,
\end{equation*}
%
то её индексы Миллера пропорциональны обратным величинам:
%
\begin{equation*}
    h:k:l
    =
    \frac{1}{n_1}
    :
    \frac{1}{n_2}
    :
    \frac{1}{n_3}.
\end{equation*}

Или, что то же самое,
%
\begin{equation*}
    n_1:n_2:n_3
    =
    \frac{1}{h}
    :
    \frac{1}{k}
    :
    \frac{1}{l}.
\end{equation*}

Поэтому выражение в скобках направлено вдоль
%
\begin{equation*}
    h\mathbf b_1
    +
    k\mathbf b_2
    +
    l\mathbf b_3.
\end{equation*}

Введём
%
\begin{equation*}
    \boxed{
        \mathbf g
        =
        h\mathbf b_1
        +
        k\mathbf b_2
        +
        l\mathbf b_3
    }.
\end{equation*}

Тогда
%
\begin{equation*}
    \mathbf c_2\times\mathbf c_3
    \parallel
    \mathbf g.
\end{equation*}

Но $\mathbf c_2\times\mathbf c_3$ является нормалью
к исходной плоскости.

Следовательно,
%
\begin{equation*}
    \boxed{
        \mathbf g\perp(hkl)
    }.
\end{equation*}

Таким образом, вектор обратной решётки
%
\begin{equation*}
    \boxed{
        \mathbf g_{hkl}
        =
        h\mathbf b_1
        +
        k\mathbf b_2
        +
        l\mathbf b_3
    }
\end{equation*}
%
перпендикулярен семейству плоскостей прямой решётки
с индексами Миллера $(hkl)$.

\textbf{Комментарий.}
Здесь устанавливается геометрический смысл обратной решётки:
направление $\mathbf g_{hkl}$ задаёт нормаль к плоскостям
$(hkl)$.

В следующей задаче увидим, что длина этого вектора связана
с расстоянием между соседними плоскостями данного семейства.


% ============================================================
% Задача 5
% ============================================================

\section{Задача 5. Расстояние между плоскостями}

\textbf{Условие.}
Если компоненты вектора $\mathbf g$ не имеют общего множителя,
показать, что $|\mathbf g|$ обратно пропорционален расстоянию
между плоскостями прямой решётки, перпендикулярными
$\mathbf g$.


% ------------------------------------------------------------
% РИСУНОК 4
%
% Вставить здесь небольшой геометрический рисунок
% из рукописного конспекта.
%
% На рисунке показать:
%   1) вектор \mathbf R;
%   2) направление \mathbf g;
%   3) проекцию \mathbf R на направление \mathbf g;
%   4) при перерисовке можно добавить семейство параллельных
%      плоскостей, перпендикулярных \mathbf g.
% ------------------------------------------------------------


Пусть вектор обратной решётки имеет вид
%
\begin{equation*}
    \mathbf g
    =
    h\mathbf b_1
    +
    k\mathbf b_2
    +
    l\mathbf b_3.
\end{equation*}

Вектор произвольного узла прямой решётки равен
%
\begin{equation*}
    \mathbf R
    =
    n_1\mathbf a_1
    +
    n_2\mathbf a_2
    +
    n_3\mathbf a_3,
    \qquad
    n_i\in\mathbb Z.
\end{equation*}

Рассмотрим их скалярное произведение:
%
\begin{align*}
    \mathbf g\cdot\mathbf R
    &=
    \left(
        h\mathbf b_1
        +
        k\mathbf b_2
        +
        l\mathbf b_3
    \right)
    \cdot
    \left(
        n_1\mathbf a_1
        +
        n_2\mathbf a_2
        +
        n_3\mathbf a_3
    \right).
\end{align*}

Полностью раскрывая скобки,
%
\begin{align*}
    \mathbf g\cdot\mathbf R
    ={}&
    hn_1\,\mathbf b_1\cdot\mathbf a_1
    +
    hn_2\,\mathbf b_1\cdot\mathbf a_2
    +
    hn_3\,\mathbf b_1\cdot\mathbf a_3
    \\
    &+
    kn_1\,\mathbf b_2\cdot\mathbf a_1
    +
    kn_2\,\mathbf b_2\cdot\mathbf a_2
    +
    kn_3\,\mathbf b_2\cdot\mathbf a_3
    \\
    &+
    ln_1\,\mathbf b_3\cdot\mathbf a_1
    +
    ln_2\,\mathbf b_3\cdot\mathbf a_2
    +
    ln_3\,\mathbf b_3\cdot\mathbf a_3.
\end{align*}

Используем основное свойство прямого и обратного базисов:
%
\begin{equation*}
    \mathbf a_i\cdot\mathbf b_j
    =
    2\pi\delta_{ij}.
\end{equation*}

Поэтому все произведения с $i\neq j$ исчезают:
%
\begin{equation*}
    \mathbf b_1\cdot\mathbf a_2
    =
    \mathbf b_1\cdot\mathbf a_3
    =
    0,
\end{equation*}
%
и аналогично для остальных векторов.

Остаётся
%
\begin{equation*}
    \mathbf g\cdot\mathbf R
    =
    2\pi hn_1
    +
    2\pi kn_2
    +
    2\pi ln_3.
\end{equation*}

Следовательно,
%
\begin{equation*}
    \boxed{
        \mathbf g\cdot\mathbf R
        =
        2\pi
        \left(
            hn_1+kn_2+ln_3
        \right)
    }.
\end{equation*}

Так как $h,k,l,n_1,n_2,n_3$ --- целые числа, величина
%
\begin{equation*}
    N
    =
    hn_1+kn_2+ln_3
\end{equation*}
%
также является целой:
%
\begin{equation*}
    N\in\mathbb Z.
\end{equation*}

Получаем
%
\begin{equation*}
    \boxed{
        \mathbf g\cdot\mathbf R
        =
        2\pi N
    }.
\end{equation*}


% ------------------------------------------------------------
% Геометрический смысл
% ------------------------------------------------------------

\subsection*{Геометрический смысл}

Уравнение
%
\begin{equation*}
    \mathbf g\cdot\mathbf r
    =
    \mathrm{const}
\end{equation*}
%
задаёт плоскость с нормалью $\mathbf g$.

Следовательно,
%
\begin{equation*}
    \boxed{
        \mathbf g\cdot\mathbf r
        =
        2\pi N
    }
\end{equation*}
%
задаёт семейство параллельных плоскостей,
перпендикулярных $\mathbf g$.

Найдём расстояние от начала координат до плоскости
с номером $N$.

Разложим $\mathbf R$ на составляющую вдоль $\mathbf g$
и составляющую, лежащую в плоскости.

В скалярное произведение входит только проекция
$\mathbf R$ на направление $\mathbf g$, поэтому
%
\begin{equation*}
    \mathbf g\cdot\mathbf R
    =
    |\mathbf g|\,d_N,
\end{equation*}
%
где $d_N$ --- расстояние от начала координат
до $N$-й плоскости.

С другой стороны,
%
\begin{equation*}
    \mathbf g\cdot\mathbf R
    =
    2\pi N.
\end{equation*}

Следовательно,
%
\begin{equation*}
    |\mathbf g|d_N
    =
    2\pi N.
\end{equation*}

Отсюда
%
\begin{equation*}
    \boxed{
        d_N
        =
        \frac{2\pi N}{|\mathbf g|}
    }.
\end{equation*}


% ------------------------------------------------------------
% Соседние плоскости
% ------------------------------------------------------------

\subsection*{Расстояние между соседними плоскостями}

Теперь найдём расстояние именно между соседними плоскостями.

Для соседней плоскости
%
\begin{equation*}
    N\longrightarrow N+1.
\end{equation*}

Поэтому
%
\begin{equation*}
    d
    =
    d_{N+1}-d_N.
\end{equation*}

Подставляем:
%
\begin{align*}
    d
    &=
    \frac{2\pi(N+1)}{|\mathbf g|}
    -
    \frac{2\pi N}{|\mathbf g|}
    \\
    &=
    \frac{2\pi}{|\mathbf g|}.
\end{align*}

Получаем
%
\begin{equation*}
    \boxed{
        d
        =
        \frac{2\pi}{|\mathbf g|}
    }.
\end{equation*}

Или
%
\begin{equation*}
    \boxed{
        |\mathbf g|
        =
        \frac{2\pi}{d}
    }.
\end{equation*}

Следовательно,
%
\begin{equation*}
    \boxed{
        |\mathbf g|
        \propto
        \frac{1}{d}
    }.
\end{equation*}

\textbf{Комментарий.}
Условие о том, что компоненты $h,k,l$ не имеют общего
множителя, здесь существенно.

Например, если
%
\begin{equation*}
    h=2h',
    \qquad
    k=2k',
    \qquad
    l=2l',
\end{equation*}
%
то выбранный вектор $\mathbf g$ соответствовал бы не ближайшим
соседним плоскостям, а каждой второй плоскости этого семейства.

Если же $h,k,l$ взаимно просты, то переход
%
\begin{equation*}
    N\longrightarrow N+1
\end{equation*}
%
действительно переводит нас к ближайшей соседней плоскости.

Поэтому для примитивного вектора данного семейства плоскостей
%
\begin{equation*}
    \boxed{
        d_{hkl}
        =
        \frac{2\pi}{
            |\mathbf g_{hkl}|
        }
    }.
\end{equation*}


% ============================================================
% Семинар 2
% ============================================================

\practice{2}{}{Колебания решётки}


% ============================================================
% Задача 1
% ============================================================

\section{Задача 1. Линейная цепочка с одним атомом в элементарной ячейке}

\textbf{Условие.}
Определить собственные моды, то есть фононный спектр,
бесконечной одномерной цепочки, состоящей из одинаковых
атомов массы $m$, соединённых одинаковыми пружинками
жёсткости $\alpha$.


% ------------------------------------------------------------
% РИСУНОК 5
%
% Вставить здесь рисунок линейной цепочки из рукописного
% конспекта / рисунок к задаче 4.1 из сборника.
%
% На рисунке:
%   1) бесконечная одномерная цепочка;
%   2) одинаковые атомы массы m;
%   3) пружины жёсткости \alpha;
%   4) расстояние между соседними атомами a;
%   5) смещение n-го атома u_n.
% ------------------------------------------------------------


Рассмотрим одномерную цепочку одинаковых атомов.

Положение равновесия $n$-го атома можно записать как
%
\begin{equation*}
    x_n=na,
\end{equation*}
%
где $a$ --- период решётки.

Пусть
%
\begin{equation*}
    u_n(t)
\end{equation*}
%
--- смещение $n$-го атома от положения равновесия.


% ------------------------------------------------------------
% Лагранжиан
% ------------------------------------------------------------

\subsection*{Лагранжиан цепочки}

Кинетическая энергия цепочки равна
%
\begin{equation*}
    T
    =
    \sum_n
    \frac{m}{2}
    \dot u_n^{\,2}.
\end{equation*}

При малых колебаниях каждую связь между соседними атомами
считаем гармонической.

Удлинение пружины между атомами $n$ и $n+1$ равно
%
\begin{equation*}
    u_{n+1}-u_n.
\end{equation*}

Поэтому потенциальная энергия этой пружины
%
\begin{equation*}
    \frac{\alpha}{2}
    (u_{n+1}-u_n)^2.
\end{equation*}

Суммарная потенциальная энергия цепочки имеет вид
%
\begin{equation*}
    U
    =
    \sum_n
    \frac{\alpha}{2}
    (u_n-u_{n+1})^2.
\end{equation*}

Следовательно,
%
\begin{equation*}
    L=T-U,
\end{equation*}
%
то есть
%
\begin{equation*}
    \boxed{
        L
        =
        \sum_n
        \left[
            \frac{m}{2}\dot u_n^{\,2}
            -
            \frac{\alpha}{2}
            (u_n-u_{n+1})^2
        \right]
    }.
\end{equation*}


% ------------------------------------------------------------
% Уравнение движения
% ------------------------------------------------------------

\subsection*{Уравнение движения}

Для каждой координаты $u_n$ записываем
уравнение Эйлера--Лагранжа:
%
\begin{equation*}
    \frac{d}{dt}
    \frac{\partial L}{\partial\dot u_n}
    -
    \frac{\partial L}{\partial u_n}
    =
    0.
\end{equation*}

Сначала вычислим производную по скорости:
%
\begin{equation*}
    \frac{\partial L}{\partial\dot u_n}
    =
    m\dot u_n.
\end{equation*}

Следовательно,
%
\begin{equation*}
    \frac{d}{dt}
    \frac{\partial L}{\partial\dot u_n}
    =
    m\ddot u_n.
\end{equation*}

Теперь найдём
%
\begin{equation*}
    \frac{\partial L}{\partial u_n}.
\end{equation*}

Здесь важно учесть, что координата $u_n$ входит
не в один, а сразу в два соседних упругих члена:
%
\begin{equation*}
    -\frac{\alpha}{2}
    (u_n-u_{n+1})^2
\end{equation*}
%
и
%
\begin{equation*}
    -\frac{\alpha}{2}
    (u_{n-1}-u_n)^2.
\end{equation*}

Поэтому
%
\begin{equation*}
    \frac{\partial L}{\partial u_n}
    =
    -\alpha(u_n-u_{n+1})
    +
    \alpha(u_{n-1}-u_n).
\end{equation*}

Раскрываем скобки:
%
\begin{align*}
    \frac{\partial L}{\partial u_n}
    &=
    -\alpha u_n
    +
    \alpha u_{n+1}
    +
    \alpha u_{n-1}
    -
    \alpha u_n
    \\
    &=
    \alpha
    \left(
        u_{n+1}
        +
        u_{n-1}
        -
        2u_n
    \right).
\end{align*}

Следовательно,
%
\begin{equation*}
    \boxed{
        \frac{\partial L}{\partial u_n}
        =
        \alpha
        \left(
            u_{n+1}
            +
            u_{n-1}
            -
            2u_n
        \right)
    }.
\end{equation*}

Получаем уравнение движения
%
\begin{equation*}
    \boxed{
        m\ddot u_n
        =
        \alpha
        \left(
            u_{n+1}
            +
            u_{n-1}
            -
            2u_n
        \right)
    }.
\end{equation*}


% ------------------------------------------------------------
% Физический смысл силы
% ------------------------------------------------------------

\subsection*{Сила, действующая на атом}

Физически правая часть уравнения движения представляет собой
результирующую силу от двух соседних пружин.

Сила от правой пружины равна
%
\begin{equation*}
    F_{n,n+1}
    =
    \alpha
    (u_{n+1}-u_n),
\end{equation*}
%
а от левой
%
\begin{equation*}
    F_{n,n-1}
    =
    \alpha
    (u_{n-1}-u_n).
\end{equation*}

Следовательно,
%
\begin{align*}
    F_n
    &=
    F_{n,n+1}
    +
    F_{n,n-1}
    \\
    &=
    \alpha
    (u_{n+1}-u_n)
    +
    \alpha
    (u_{n-1}-u_n)
    \\
    &=
    \alpha
    \left(
        u_{n+1}
        +
        u_{n-1}
        -
        2u_n
    \right).
\end{align*}

Таким образом, снова получаем
%
\begin{equation*}
    m\ddot u_n
    =
    \alpha
    \left(
        u_{n+1}
        +
        u_{n-1}
        -
        2u_n
    \right).
\end{equation*}


% ------------------------------------------------------------
% Нормальная мода
% ------------------------------------------------------------

\subsection*{Нормальная мода}

Ищем нормальное колебание в виде плоской волны:
%
\begin{equation*}
    \boxed{
        u_n
        =
        u
        e^{-i\omega t+ikna}
    }.
\end{equation*}

Здесь $k$ --- волновой вектор колебания,
$\omega$ --- частота нормальной моды,
а $a$ --- период цепочки.

\textbf{Комментарий.}
В представленном рукописном конспекте решение заканчивается
именно на этом месте --- после записи плоской волны
%
\begin{equation*}
    u_n
    =
    u
    e^{-i\omega t+ikna}.
\end{equation*}

Дальнейшая подстановка этого выражения в уравнение движения
и получение дисперсионного соотношения $\omega(k)$ на
приложенной странице уже не записаны.


\end{document}